\documentclass[11pt]{article}
\usepackage[T1]{fontenc}
\usepackage{amsmath,amssymb}
\usepackage[margin=1in]{geometry}
\usepackage[hidelinks]{hyperref}
\hypersetup{pdftitle={Development summary: stick-knot counting},
 pdfauthor={},pdfsubject={Retrospective account based on project records}}
\ifdefined\pdftrailerid\pdftrailerid{}\fi
\setlength{\emergencystretch}{2em}
\title{Development summary: stick-knot counting}
\author{}
\date{23 October 2026}
\begin{document}
\maketitle
\noindent\textbf{Manuscript authors: Alexander Kolpakov and Igor Rivin.}

\noindent\textbf{Discussions and model prompting: Alexander Kolpakov (UATX).}
This retrospective account summarizes the released manuscript and its
supporting mathematical documentation. It is not a transcript or a reconstruction of model
turns. No original conversation export was available in the source tree.

\section*{Schematic research brief}
The task is to determine the growth of the number of ordinary knot types
realizable with a budget of at most $N$ straight sticks. A schematic research
brief is to find an upper bound from polygon geometry, a complementary lower
construction, and a precise account of the mathematical inputs. This is an
editorial description of the task, not a quotation of an original prompt.

For the set $\mathcal K_N$ of these knot types, the manuscript presents
\[
 N^{(2/3+o(1))N}\le |\mathcal K_N|\le N^{3N+o(N)},
 \qquad \log |\mathcal K_N|=\Theta(N\log N).
\]
This determines the logarithmic growth order. It does not determine a leading
constant or an asymptotic equivalent for the number of knot types.

\section*{Mathematical development}
\begin{enumerate}
\item \textbf{Turn knot geometry into a low-degree wall problem.}
  Describe an $N$-vertex polygon by $3N$ real coordinates. Cubic determinants
  detect the possible intersection of nonincident edges, while quartic
  squared cross products exclude collapsed or overlapping adjacent edges.
  The number of walls grows only quadratically in $N$.
\item \textbf{Pass from chambers to knot types.}
  Every bounded-stick knot has a generic representative in the wall-free
  open locus. Knot type is constant along paths, so a surjection from
  connected components gives an upper bound for the knot population.
\item \textbf{Keep the factorial saving in the component count.}
  Apply the refined real-algebraic component estimate with its partial
  binomial sum retained. The dimension factorial changes the relevant
  base from quadratic to linear in $N$, producing the upper exponent 3.
\item \textbf{Build a large family of distinct braid words.}
  Four independent permutations on $k$ letters give at least $(k!)^4$
  doubly alternating permutations on $4k$ letters. Starting and finishing
  conditions make the corresponding words distinct normal forms.
\item \textbf{Control the endpoint and stick costs.}
  Pigeonhole by endpoint permutation and append a common correcting factor
  to obtain distinct pure braids with controlled length. Transfer these to
  knots through the cited braid-to-knot and arc-to-stick results, including
  the bounded loss on forgetting orientation.
\item \textbf{Optimize and check the counting deductions.}
  Suitable fixed word lengths give every lower coefficient below $2/3$.
  A logarithmically varying word length yields the refined lower deficit
  $-(2/3)N\log\log N-O(N)$. The resulting upper and lower estimates give
  the common $N\log N$ scale for the logarithm.
\end{enumerate}

The upper argument counts discriminant chambers directly. It is independent
of related strong-geometry reconstruction results. The manuscript states
its priority claim with the qualification ``to our knowledge'' and keeps
the original real-algebraic, braid, and knot-theoretic inputs attributed.

\section*{Verification and scope}
The standalone Lean project uses actual polygon coordinates, wall
polynomials, paths and components, permutation words, endpoint fibers, and
a bounded-stick subtype of a supplied knot interface. It checks the internal
finite and asymptotic deductions, including the optimized floor/log estimate.

Ambient-isotopy invariance, the Barone--Basu component theorem, structural
Garside facts, the pure-braid-to-knot injection, the arc-to-stick inequality,
and orientation reversal enter through explicit external mathematical
hypotheses. The classical ambient-isotopy quotient and the auxiliary link
construction are outside the formal target. The formal certificate is
conditional on those stated inputs.

The manuscript authors are Alexander Kolpakov and Igor Rivin. Their existing
acknowledgment reports GPT-5.6-sol assistance as a research and coding tool.
This summary describes released project records; it does not authenticate
an original model conversation. Account and session metadata, credentials,
local paths, detailed repair logs, and operational information are omitted.

\end{document}
